\chapter*{Abstract}
\addcontentsline{toc}{chapter}{Abstract}

This report presents \textbf{Cognitive Copilot}, a full-stack academic learning platform that combines core Learning Management System (LMS) workflows with an AI tutor. The project is built around practical university needs: course and topic management, classroom community interactions, routine and attendance support, notifications, and AI-assisted study. The system is implemented as a web application with a React frontend and a FastAPI backend.

The implementation follows a modular architecture. The frontend provides role-based pages for students and tutors, while the backend exposes REST APIs for authentication, courses, routine, community, analytics, settings, profile, materials, and notifications. For AI functionality, the platform integrates local model inference with retrieval-based question answering and quiz support in the AI Tutor workspace. OCR-assisted ingestion and asynchronous workers support long-running operations so user interactions stay responsive.

The project emphasizes engineering reliability over experimental claims. The report therefore focuses on implemented modules, data flow, CRUD behaviors, validation and authorization design, and tested user workflows. No unsupported performance or accuracy numbers are presented as final evidence; instead, verification is described through reproducible build and test procedures across backend and frontend.

Cognitive Copilot demonstrates how a project can remain academically meaningful while also being implementation-driven: it delivers a unified learning environment, keeps AI features integrated with course context, and preserves maintainable software structure suitable for further extension in future academic work.
